% Options for packages loaded elsewhere
\PassOptionsToPackage{unicode}{hyperref}
\PassOptionsToPackage{hyphens}{url}
\PassOptionsToPackage{dvipsnames,svgnames,x11names}{xcolor}
%
\documentclass[
  12pt,
  letterpaper]{letter}

\usepackage{amsmath,amssymb}
\usepackage{iftex}
\ifPDFTeX
  \usepackage[T1]{fontenc}
  \usepackage[utf8]{inputenc}
  \usepackage{textcomp} % provide euro and other symbols
\else % if luatex or xetex
  \usepackage{unicode-math}
  \defaultfontfeatures{Scale=MatchLowercase}
  \defaultfontfeatures[\rmfamily]{Ligatures=TeX,Scale=1}
\fi
\usepackage{lmodern}
\ifPDFTeX\else  
    % xetex/luatex font selection
\fi
% Use upquote if available, for straight quotes in verbatim environments
\IfFileExists{upquote.sty}{\usepackage{upquote}}{}
\IfFileExists{microtype.sty}{% use microtype if available
  \usepackage[]{microtype}
  \UseMicrotypeSet[protrusion]{basicmath} % disable protrusion for tt fonts
}{}
\makeatletter
\@ifundefined{KOMAClassName}{% if non-KOMA class
  \IfFileExists{parskip.sty}{%
    \usepackage{parskip}
  }{% else
    \setlength{\parindent}{0pt}
    \setlength{\parskip}{6pt plus 2pt minus 1pt}}
}{% if KOMA class
  \KOMAoptions{parskip=half}}
\makeatother
\usepackage{xcolor}
\usepackage[margin=1in,bottom=1in,left=1in,right=1in]{geometry}
\setlength{\emergencystretch}{3em} % prevent overfull lines
\setcounter{secnumdepth}{-\maxdimen} % remove section numbering
% Make \paragraph and \subparagraph free-standing
\makeatletter
\ifx\paragraph\undefined\else
  \let\oldparagraph\paragraph
  \renewcommand{\paragraph}{
    \@ifstar
      \xxxParagraphStar
      \xxxParagraphNoStar
  }
  \newcommand{\xxxParagraphStar}[1]{\oldparagraph*{#1}\mbox{}}
  \newcommand{\xxxParagraphNoStar}[1]{\oldparagraph{#1}\mbox{}}
\fi
\ifx\subparagraph\undefined\else
  \let\oldsubparagraph\subparagraph
  \renewcommand{\subparagraph}{
    \@ifstar
      \xxxSubParagraphStar
      \xxxSubParagraphNoStar
  }
  \newcommand{\xxxSubParagraphStar}[1]{\oldsubparagraph*{#1}\mbox{}}
  \newcommand{\xxxSubParagraphNoStar}[1]{\oldsubparagraph{#1}\mbox{}}
\fi
\makeatother


\providecommand{\tightlist}{%
  \setlength{\itemsep}{0pt}\setlength{\parskip}{0pt}}\usepackage{longtable,booktabs,array}
\usepackage{calc} % for calculating minipage widths
% Correct order of tables after \paragraph or \subparagraph
\usepackage{etoolbox}
\makeatletter
\patchcmd\longtable{\par}{\if@noskipsec\mbox{}\fi\par}{}{}
\makeatother
% Allow footnotes in longtable head/foot
\IfFileExists{footnotehyper.sty}{\usepackage{footnotehyper}}{\usepackage{footnote}}
\makesavenoteenv{longtable}
\usepackage{graphicx}
\makeatletter
\def\maxwidth{\ifdim\Gin@nat@width>\linewidth\linewidth\else\Gin@nat@width\fi}
\def\maxheight{\ifdim\Gin@nat@height>\textheight\textheight\else\Gin@nat@height\fi}
\makeatother
% Scale images if necessary, so that they will not overflow the page
% margins by default, and it is still possible to overwrite the defaults
% using explicit options in \includegraphics[width, height, ...]{}
\setkeys{Gin}{width=\maxwidth,height=\maxheight,keepaspectratio}
% Set default figure placement to htbp
\makeatletter
\def\fps@figure{htbp}
\makeatother

\usepackage{eso-pic} % for background image
\usepackage{graphicx} % necessary for images
\usepackage{xifthen}
\usepackage{parskip} % for spacing
\usepackage{fancyhdr} % for fancy headers
\usepackage{lastpage} 
\usepackage{bera} % for font
\pagestyle{fancy}
\renewcommand{\headrulewidth}{0pt}
\fancyhead{} % footer for pages 2 on
\lfoot{} % Left footer
\cfoot{} % Change for center footer
\rfoot{\footnotesize Page \thepage\ of \pageref{LastPage}} % Right footer page #s 
\renewcommand{\footrulewidth}{0pt}
\makeatletter
\@ifpackageloaded{caption}{}{\usepackage{caption}}
\AtBeginDocument{%
\ifdefined\contentsname
  \renewcommand*\contentsname{Table of contents}
\else
  \newcommand\contentsname{Table of contents}
\fi
\ifdefined\listfigurename
  \renewcommand*\listfigurename{List of Figures}
\else
  \newcommand\listfigurename{List of Figures}
\fi
\ifdefined\listtablename
  \renewcommand*\listtablename{List of Tables}
\else
  \newcommand\listtablename{List of Tables}
\fi
\ifdefined\figurename
  \renewcommand*\figurename{Figure}
\else
  \newcommand\figurename{Figure}
\fi
\ifdefined\tablename
  \renewcommand*\tablename{Table}
\else
  \newcommand\tablename{Table}
\fi
}
\@ifpackageloaded{float}{}{\usepackage{float}}
\floatstyle{ruled}
\@ifundefined{c@chapter}{\newfloat{codelisting}{h}{lop}}{\newfloat{codelisting}{h}{lop}[chapter]}
\floatname{codelisting}{Listing}
\newcommand*\listoflistings{\listof{codelisting}{List of Listings}}
\makeatother
\makeatletter
\makeatother
\makeatletter
\@ifpackageloaded{caption}{}{\usepackage{caption}}
\@ifpackageloaded{subcaption}{}{\usepackage{subcaption}}
\makeatother

\ifLuaTeX
  \usepackage{selnolig}  % disable illegal ligatures
\fi
\usepackage{bookmark}

\IfFileExists{xurl.sty}{\usepackage{xurl}}{} % add URL line breaks if available
\urlstyle{same} % disable monospaced font for URLs
\hypersetup{
  pdftitle={Letter from ANU},
  pdfauthor={\hspace{0pt}},
  colorlinks=true,
  linkcolor={blue},
  filecolor={Maroon},
  citecolor={Blue},
  urlcolor={Blue},
  pdfcreator={LaTeX via pandoc}}


\title{Letter from ANU}
\author{\hspace{0pt}}
\date{July 29, 2025}

\begin{document}
\AddToShipoutPicture{\ifthenelse{\value{page}<2}{\includegraphics[width=\paperwidth,height=\paperheight]{\_extensions/quarto-anu-letter/assets/images/coverpage.jpg}}{\includegraphics[width=\paperwidth,height=\paperheight]{\_extensions/quarto-anu-letter/assets/images/page-background.jpg}}}


\signature{
\hspace{0pt}}
\address{Australian National University\\Research School of Finance,
Actuarial Studies and Statistics\\26C Kingsley
Street\\Canberra 2600\\Australia\\[2mm]}
\begin{letter}{Editors,\\Journal of Computational and Graphical
Statistics}
\opening{Dear Prof.~Chen and Prof.~Sangalli,}


Please consider our manuscript ``Automated Assessment of Residual Plots
with Computer Vision Models'' for publication in the Journal of
Computational and Graphical Statistics.

We confirm that this work is original and has not been published
elsewhere, nor is it currently under consideration for publication
elsewhere.

This paper describes a new computer vision model that can automatically
read residual plots, it's architecture and performance relative to human
evaluation and classical hypothesis testing. Our previous work, ``A Plot
is Worth a Thousand Tests: Assessing Residual Diagnostics with the
Lineup Protocol'' (JCGS, 2024), described interesting findings about
human evaluations of residual plots relative to classical hypothesis
testing, discovered when building training samples for the computer
vision model described in our new paper. The computer vision model
described in this new submission was our primary research goal.

We specifically trained a custom model to predict a specified distance
measure that quantifies the discrepancy between residuals and
theoretically ``good'' distributions. Building on these predictions, we
construct a statistical testing framework aligned with the original
lineup protocol, ensuring a valid visual test. Despite substantial
experimentation, we could not manage to get a computer vision model to
perform quite as well as the human evaluations. It is slightly more
sensitive, flagging more residual plots as bad. However, as a computer
model, it can process substantially more residual plots than is possible
through human evaluation. Relative to statistical tests (Breusch-Pagan,
Ramsey RESET, Shapiro-Wilk) it performs substantially better, is less
sensitive to small departures, and jointly assesses a variety of
departures (heteroscedasticity, non-linearity, non-normality).

The paper also shows that the model, like human observers, relies on
visual patterns and shapes, making it a practical tool for reducing
manual effort in residual plot diagnostics. This approach would be of
particular interest to readers concerned with model diagnostics and the
automation of visual data analysis. It represents a new direction for
statistical graphics research that engages with new computational tools.
There is a shiny app and R package available to support access and use
of the computer vision model, and this is documented in a special issue
of the Australian and New Zealand Journal of Statistics that celebrates
25 years of R, to appear later this year.

We have no conflict of interest to disclose.

Please address all corresponding concerning this manuscript to me at
patrick.li@anu.edu.au.

Thank you for the consideration of this manuscript. We believe that it
is a good fit for the Journal of Computational and Graphical Statistics.



\closing{Sincerely,}
\vfill
\end{letter}
\end{document}
